% ============================================================
% SemArq1_TeamBales.tex — Documento principal
% Clase: report
%
% Orden exigido por la orientación del seminario:
%   portada → índice → introducción → preguntas P1-P6 → conclusiones
%   → referencias bibliográficas → Declaración de Uso de la IA
%   (la declaración va SIEMPRE después de la bibliografía).
%
% Cada pregunta se incluye solo si su solucion.tex ya está en el
% repositorio: así main compila aunque falten capítulos, y cada PR de
% pregunta los añade sin tocar este archivo. \pregunta fija además el
% número del capítulo, para que P4 sea el capítulo 4 aunque falten los
% anteriores. Se usa \input (no \include) para que los capítulos vayan
% seguidos, sin salto de página.
% ============================================================

\documentclass[12pt]{report}

% ============================================================
% preamble.tex — Preámbulo compartido de SemArq1_TeamBales
% Los fragmentos (solucion.tex, notas.tex) NO cargan paquetes:
% todo lo que necesiten se carga aquí.
% ============================================================

% --- Página: márgenes de 2,5 cm (la portada alinea su texto con este margen) ---
\usepackage[margin=2.5cm]{geometry}

% --- Codificación y fuentes (ATS: tildes extraíbles con pdftotext) ---
\usepackage[T1]{fontenc}
\usepackage{lmodern}

% --- Idioma ---
% es-nolists: babel no redefine las listas (enumitem las controla).
% es-tabla:   «Tabla» en lugar de «Cuadro».
\usepackage[spanish, es-nolists, es-tabla]{babel}
\usepackage{csquotes}
\addto\captionsspanish{\renewcommand{\contentsname}{Índice}}

% --- Bibliografía: autor-año, una sola lista final ---
\usepackage[backend=biber, style=authoryear, natbib=true]{biblatex}
\addbibresource{core.bib}
% Fuentes nuevas de cada pregunta e investigación (solo si el archivo existe).
\newcommand{\bibsiexiste}[1]{\IfFileExists{#1}{\addbibresource{#1}}{}}
\bibsiexiste{preguntas/p1-trade-offs/referencias.bib}
\bibsiexiste{preguntas/p2-simetria-hexagono/referencias.bib}
\bibsiexiste{preguntas/p3-ddd-bounded-context/referencias.bib}
\bibsiexiste{preguntas/p4-inversion-dependencias/referencias.bib}
\bibsiexiste{preguntas/p5-kitty-problem/referencias.bib}
\bibsiexiste{preguntas/p6-ncapas-calidad/referencias.bib}
\bibsiexiste{investigacion/p1-regla-dependencia/fuentes.bib}
\bibsiexiste{investigacion/p2-hexagonal-onion/fuentes.bib}
\bibsiexiste{investigacion/p3-bounded-context/fuentes.bib}
\bibsiexiste{investigacion/p4-dip-interfaces/fuentes.bib}
\bibsiexiste{investigacion/p5-kitty-ocp/fuentes.bib}
\bibsiexiste{investigacion/p6-open-layers/fuentes.bib}

% --- Matemáticas, tablas, figuras y código ---
\usepackage{amsmath}
\usepackage{graphicx}
\usepackage{booktabs}
\usepackage{array}
\usepackage{tabularx}
\usepackage{float}
\usepackage{enumitem}
\usepackage{listings}
\usepackage{xcolor}
\usepackage{tikz}
% babel: evita el choque entre los atajos de babel-spanish (< y >) y las flechas de TikZ.
\usetikzlibrary{arrows.meta, positioning, shapes.geometric, calc, babel}

% --- Microtipografía: mejora el ajuste de las líneas (texto real, sin efecto en ATS) ---
\usepackage{microtype}

% --- Hipervínculos (al final, como recomienda hyperref) ---
\usepackage{hyperref}

% Estilos compartidos de diagramas (copiables por P3 y P4)
\tikzset{
  bloque/.style={rectangle, draw, thick, minimum width=2.4cm, minimum height=1.1cm, align=center, font=\sffamily\small},
  flecha/.style={-{Stealth[length=3mm]}, thick, shorten >=1pt, shorten <=1pt},
  capa/.style={rectangle, draw=gray!60, thick, rounded corners, minimum width=3.2cm, minimum height=1.6cm, align=center, font=\sffamily\footnotesize, fill=gray!8},
  contexto/.style={draw, thick, rounded corners=4pt, minimum width=4cm, minimum height=2.6cm, align=center, font=\sffamily\small, fill=blue!6},
  linea-integracion/.style={-{Stealth}, thick, dashed},
  acl/.style={rectangle, draw, thick, minimum width=2.6cm, minimum height=1cm, align=center, font=\sffamily\footnotesize, fill=orange!15},
  dependencia/.style={-{Latex[length=2.5mm]}, thick, shorten >=1pt, shorten <=1pt},
  modulo/.style={rectangle, draw, thick, minimum width=2.8cm, minimum height=0.9cm, align=center, font=\sffamily\small},
  interfaz/.style={rectangle, draw, thick, dashed, minimum width=2.8cm, minimum height=0.9cm, align=center, font=\sffamily\small\itshape, fill=yellow!12},
}

% Ajustes de composición
% - \emergencystretch: permite repartir más espacio en líneas con palabras que no se
%   pueden cortar (p. ej. nombres de clases en \texttt), en vez de salirse del margen.
% - biburl*penalty: permite cortar las URL largas de la bibliografía.
\setlength{\emergencystretch}{3em}
% - \tolerance: acepta líneas algo más espaciadas antes que dejar una que se salga del margen.
\tolerance=1000
\setcounter{biburlnumpenalty}{100}
\setcounter{biburlucpenalty}{100}
\setcounter{biburllcpenalty}{100}

% Capítulos seguidos, sin salto de página: cada capítulo (introducción, preguntas,
% conclusiones, bibliografía, uso de IA) empieza a continuación del anterior.
% Se quita el \clearpage de \chapter y se reduce el espacio antes del título.
\makeatletter
\patchcmd{\chapter}{\if@openright\cleardoublepage\else\clearpage\fi}{\par\bigskip}{}{}
\patchcmd{\@makeschapterhead}{\vspace*{50\p@}}{\vspace*{10\p@}}{}{}
% Títulos de capítulo: sin la línea «Capítulo N» (el título ya dice «P4: …»),
% en \LARGE en lugar de \Huge y con menos espacio. El número del capítulo se
% conserva internamente para numerar las secciones (4.1, 4.2…) y el índice.
\renewcommand{\@makechapterhead}[1]{%
  \vspace*{10\p@}%
  {\parindent \z@ \raggedright \normalfont
    \interlinepenalty\@M
    \LARGE\bfseries #1\par\nobreak
    \vskip 20\p@}}
\patchcmd{\@makeschapterhead}{\Huge \bfseries #1}{\LARGE\bfseries #1}{}{}
\patchcmd{\@makeschapterhead}{\vskip 40\p@}{\vskip 20\p@}{}{}
\makeatother

% Formato único de citas: (Autor, año, p. XX) y (Autor, año, pp. XX–YY).
% - Prefijos «p.» y «pp.» también cuando biblatex los añade solo (en español pondría
%   «pág.» y «págs.»).
% - \mknormrange: los rangos de páginas se escriben con guion largo (154–155).
\DefineBibliographyStrings{spanish}{page = {p\adddot}, pages = {pp\adddot}}
% biblatex en español separa los rangos con guion corto; aquí se usa el guion largo.
\DefineBibliographyExtras{spanish}{\protected\def\bibrangedash{\textendash\penalty\hyphenpenalty}}
\DeclareFieldFormat{postnote}{\mkpageprefix[pagination][\mknormrange]{#1}}
\DeclareFieldFormat{multipostnote}{\mkpageprefix[pagination][\mknormrange]{#1}}

% Wrappers de citas: biblatex elige «p.» (una página) o «pp.» (rango o lista).
\newcommand{\fuente}[2]{\citep[#2]{#1}}
\newcommand{\nombrefuente}[1]{\citet{#1}}

% Referencias bibliográficas en formato APA con viñetas
\defbibenvironment{bibliography}
  {\list
     {\textbullet}
     {\setlength{\labelwidth}{1em}%
      \setlength{\labelsep}{0.5em}%
      \setlength{\leftmargin}{1.5em}%
      \setlength{\itemindent}{0pt}%
      \setlength{\itemsep}{0.4em}%
      \setlength{\parsep}{0pt}}}
  {\endlist}
  {\item}

% Metadatos PDF (ATS)
\hypersetup{
  colorlinks=true,
  linkcolor=blue,
  citecolor=black,
  urlcolor=blue,
  pdftitle={Seminario de Arquitecturas de Software: Un Análisis Comparativo de Estilos Basados en Capas y Dominio},
  pdfauthor={Team Bales},
}

% ============================================================
% portada.tex — Página de presentación de SemArq1_TeamBales
% Propiedad del issue #13. Se carga con % ============================================================
% portada.tex — Página de presentación de SemArq1_TeamBales
% Propiedad del issue #13. Se carga con % ============================================================
% portada.tex — Página de presentación de SemArq1_TeamBales
% Propiedad del issue #13. Se carga con \input{portada} en el
% preámbulo y se usa con \portada tras \begin{document}.
%
% Cumple los 5 puntos de «Página de presentación» de la
% orientación del seminario (sección 3):
%   1. Nombre de la actividad      4. Contribución de cada integrante
%   2. Nombre del equipo           5. Enlace al repositorio público
%   3. Nombres completos de los integrantes
% Todo el texto es real (no imágenes): apto para ATS.
% ============================================================

\definecolor{portadaazul}{HTML}{1F3A5F}
\definecolor{portadaazulclaro}{HTML}{3B5F8A}
\definecolor{portadaacento}{HTML}{E0A526}
\definecolor{portadagris}{HTML}{5B6573}

\newcommand{\portada}{%
  \begin{titlepage}
    \setlength{\parindent}{0pt}%
    % --- Fondo: banda superior con encabezado y motivo, franja de acento y pie ---
    \begin{tikzpicture}[remember picture, overlay]
      \coordinate (bandaSO) at ([yshift=-5.4cm]current page.north west);
      \fill[portadaazul] (current page.north west) rectangle ([yshift=-5.4cm]current page.north east);
      % Motivo dentro de la banda: anillos (Onion) y hexágono (Hexagonal)
      \begin{scope}
        \clip (current page.north west) rectangle ([yshift=-5.4cm]current page.north east);
        \foreach \r in {0.8, 1.5, 2.2, 2.9, 3.6}
          {\draw[portadaazulclaro, line width=1.1pt] ([xshift=-2.6cm, yshift=-2.9cm]current page.north east) circle (\r);}
        \node[regular polygon, regular polygon sides=6, minimum size=3.0cm,
              draw=portadaacento, line width=1.6pt]
              at ([xshift=-6.3cm, yshift=-3.3cm]current page.north east) {};
      \end{scope}
      \fill[portadaacento] (bandaSO) rectangle ([yshift=-5.65cm]current page.north east);
      % Encabezado institucional
      \node[anchor=south west, align=left, text=white, font=\sffamily]
           at ([xshift=2.5cm, yshift=0.9cm]bandaSO)
           {{\Large\bfseries Universidad de La Habana}\\[4pt]
            {\large Facultad de Matemática y Computación}\\[10pt]
            {\normalsize Ingeniería de Software\ \ \textbar\ \ 3.\textsuperscript{er} año\ \ \textbar\ \ Curso 2026--2027}};
      % Pie
      \fill[portadaazul] (current page.south west) rectangle ([yshift=1.3cm]current page.south east);
      \node[anchor=west, text=white, font=\sffamily\small]
           at ([xshift=2.5cm, yshift=0.65cm]current page.south west) {La Habana, \today};
      \node[anchor=east, text=white, font=\sffamily\small]
           at ([xshift=-2.5cm, yshift=0.65cm]current page.south east) {SemArq1\_TeamBales};
    \end{tikzpicture}

    \vspace*{4.2cm}

    % --- 1. Nombre de la actividad ---
    {\sffamily\small\bfseries\color{portadaacento} ACTIVIDAD\par}
    \vspace{0.3cm}
    {\sffamily\fontsize{26}{31}\selectfont\bfseries\color{portadaazul}
      Seminario de Arquitecturas\\[2pt] de Software:\par}
    \vspace{0.4cm}
    {\sffamily\Large\color{portadagris}
      Un Análisis Comparativo de Estilos\\[2pt] Basados en Capas y Dominio\par}
    \vspace{0.55cm}
    {\color{portadaacento}\rule{3.2cm}{2.5pt}\par}

    \vspace{0.9cm}

    % --- 2. Nombre del equipo ---
    {\sffamily\small\color{portadagris} EQUIPO\par}
    \vspace{0.12cm}
    {\sffamily\huge\bfseries\color{portadaazul} Team Bales\par}

    \vspace{0.8cm}

    % --- 3 y 4. Integrantes y contribución individual ---
    {\sffamily\small\color{portadagris} INTEGRANTES Y PREGUNTAS EN LAS QUE TRABAJÓ\par}
    \vspace{0.3cm}
    {\sffamily\normalsize
      \newcommand{\integrante}[2]{\textbf{##1}\quad{\color{portadaacento}\textbullet}\quad{\color{portadagris}##2}\par\vspace{0.22cm}}%
      \integrante{Heily Rodríguez Rodríguez}{1 y 2}
      \integrante{Enrique González Moreira}{4 y 6}
      \integrante{Ernesto Alejandro Soler Choong}{3 y 5}
      \integrante{Mauricio Brindon Carbo}{Sin contribución}
    }

    \vspace{0.8cm}

    % --- 5. Enlace al repositorio público ---
    {\sffamily\small\color{portadagris} REPOSITORIO PÚBLICO\par}
    \vspace{0.12cm}
    {\sffamily\large\href{https://github.com/Team-Bales/seminario-arq-capas}%
      {\color{portadaazul}https://github.com/Team-Bales/seminario-arq-capas}\par}
  \end{titlepage}
}
 en el
% preámbulo y se usa con \portada tras \begin{document}.
%
% Cumple los 5 puntos de «Página de presentación» de la
% orientación del seminario (sección 3):
%   1. Nombre de la actividad      4. Contribución de cada integrante
%   2. Nombre del equipo           5. Enlace al repositorio público
%   3. Nombres completos de los integrantes
% Todo el texto es real (no imágenes): apto para ATS.
% ============================================================

\definecolor{portadaazul}{HTML}{1F3A5F}
\definecolor{portadaazulclaro}{HTML}{3B5F8A}
\definecolor{portadaacento}{HTML}{E0A526}
\definecolor{portadagris}{HTML}{5B6573}

\newcommand{\portada}{%
  \begin{titlepage}
    \setlength{\parindent}{0pt}%
    % --- Fondo: banda superior con encabezado y motivo, franja de acento y pie ---
    \begin{tikzpicture}[remember picture, overlay]
      \coordinate (bandaSO) at ([yshift=-5.4cm]current page.north west);
      \fill[portadaazul] (current page.north west) rectangle ([yshift=-5.4cm]current page.north east);
      % Motivo dentro de la banda: anillos (Onion) y hexágono (Hexagonal)
      \begin{scope}
        \clip (current page.north west) rectangle ([yshift=-5.4cm]current page.north east);
        \foreach \r in {0.8, 1.5, 2.2, 2.9, 3.6}
          {\draw[portadaazulclaro, line width=1.1pt] ([xshift=-2.6cm, yshift=-2.9cm]current page.north east) circle (\r);}
        \node[regular polygon, regular polygon sides=6, minimum size=3.0cm,
              draw=portadaacento, line width=1.6pt]
              at ([xshift=-6.3cm, yshift=-3.3cm]current page.north east) {};
      \end{scope}
      \fill[portadaacento] (bandaSO) rectangle ([yshift=-5.65cm]current page.north east);
      % Encabezado institucional
      \node[anchor=south west, align=left, text=white, font=\sffamily]
           at ([xshift=2.5cm, yshift=0.9cm]bandaSO)
           {{\Large\bfseries Universidad de La Habana}\\[4pt]
            {\large Facultad de Matemática y Computación}\\[10pt]
            {\normalsize Ingeniería de Software\ \ \textbar\ \ 3.\textsuperscript{er} año\ \ \textbar\ \ Curso 2026--2027}};
      % Pie
      \fill[portadaazul] (current page.south west) rectangle ([yshift=1.3cm]current page.south east);
      \node[anchor=west, text=white, font=\sffamily\small]
           at ([xshift=2.5cm, yshift=0.65cm]current page.south west) {La Habana, \today};
      \node[anchor=east, text=white, font=\sffamily\small]
           at ([xshift=-2.5cm, yshift=0.65cm]current page.south east) {SemArq1\_TeamBales};
    \end{tikzpicture}

    \vspace*{4.2cm}

    % --- 1. Nombre de la actividad ---
    {\sffamily\small\bfseries\color{portadaacento} ACTIVIDAD\par}
    \vspace{0.3cm}
    {\sffamily\fontsize{26}{31}\selectfont\bfseries\color{portadaazul}
      Seminario de Arquitecturas\\[2pt] de Software:\par}
    \vspace{0.4cm}
    {\sffamily\Large\color{portadagris}
      Un Análisis Comparativo de Estilos\\[2pt] Basados en Capas y Dominio\par}
    \vspace{0.55cm}
    {\color{portadaacento}\rule{3.2cm}{2.5pt}\par}

    \vspace{0.9cm}

    % --- 2. Nombre del equipo ---
    {\sffamily\small\color{portadagris} EQUIPO\par}
    \vspace{0.12cm}
    {\sffamily\huge\bfseries\color{portadaazul} Team Bales\par}

    \vspace{0.8cm}

    % --- 3 y 4. Integrantes y contribución individual ---
    {\sffamily\small\color{portadagris} INTEGRANTES Y PREGUNTAS EN LAS QUE TRABAJÓ\par}
    \vspace{0.3cm}
    {\sffamily\normalsize
      \newcommand{\integrante}[2]{\textbf{##1}\quad{\color{portadaacento}\textbullet}\quad{\color{portadagris}##2}\par\vspace{0.22cm}}%
      \integrante{Heily Rodríguez Rodríguez}{1 y 2}
      \integrante{Enrique González Moreira}{4 y 6}
      \integrante{Ernesto Alejandro Soler Choong}{3 y 5}
      \integrante{Mauricio Brindon Carbo}{Sin contribución}
    }

    \vspace{0.8cm}

    % --- 5. Enlace al repositorio público ---
    {\sffamily\small\color{portadagris} REPOSITORIO PÚBLICO\par}
    \vspace{0.12cm}
    {\sffamily\large\href{https://github.com/Team-Bales/seminario-arq-capas}%
      {\color{portadaazul}https://github.com/Team-Bales/seminario-arq-capas}\par}
  \end{titlepage}
}
 en el
% preámbulo y se usa con \portada tras \begin{document}.
%
% Cumple los 5 puntos de «Página de presentación» de la
% orientación del seminario (sección 3):
%   1. Nombre de la actividad      4. Contribución de cada integrante
%   2. Nombre del equipo           5. Enlace al repositorio público
%   3. Nombres completos de los integrantes
% Todo el texto es real (no imágenes): apto para ATS.
% ============================================================

\definecolor{portadaazul}{HTML}{1F3A5F}
\definecolor{portadaazulclaro}{HTML}{3B5F8A}
\definecolor{portadaacento}{HTML}{E0A526}
\definecolor{portadagris}{HTML}{5B6573}

\newcommand{\portada}{%
  \begin{titlepage}
    \setlength{\parindent}{0pt}%
    % --- Fondo: banda superior con encabezado y motivo, franja de acento y pie ---
    \begin{tikzpicture}[remember picture, overlay]
      \coordinate (bandaSO) at ([yshift=-5.4cm]current page.north west);
      \fill[portadaazul] (current page.north west) rectangle ([yshift=-5.4cm]current page.north east);
      % Motivo dentro de la banda: anillos (Onion) y hexágono (Hexagonal)
      \begin{scope}
        \clip (current page.north west) rectangle ([yshift=-5.4cm]current page.north east);
        \foreach \r in {0.8, 1.5, 2.2, 2.9, 3.6}
          {\draw[portadaazulclaro, line width=1.1pt] ([xshift=-2.6cm, yshift=-2.9cm]current page.north east) circle (\r);}
        \node[regular polygon, regular polygon sides=6, minimum size=3.0cm,
              draw=portadaacento, line width=1.6pt]
              at ([xshift=-6.3cm, yshift=-3.3cm]current page.north east) {};
      \end{scope}
      \fill[portadaacento] (bandaSO) rectangle ([yshift=-5.65cm]current page.north east);
      % Encabezado institucional
      \node[anchor=south west, align=left, text=white, font=\sffamily]
           at ([xshift=2.5cm, yshift=0.9cm]bandaSO)
           {{\Large\bfseries Universidad de La Habana}\\[4pt]
            {\large Facultad de Matemática y Computación}\\[10pt]
            {\normalsize Ingeniería de Software\ \ \textbar\ \ 3.\textsuperscript{er} año\ \ \textbar\ \ Curso 2026--2027}};
      % Pie
      \fill[portadaazul] (current page.south west) rectangle ([yshift=1.3cm]current page.south east);
      \node[anchor=west, text=white, font=\sffamily\small]
           at ([xshift=2.5cm, yshift=0.65cm]current page.south west) {La Habana, \today};
      \node[anchor=east, text=white, font=\sffamily\small]
           at ([xshift=-2.5cm, yshift=0.65cm]current page.south east) {SemArq1\_TeamBales};
    \end{tikzpicture}

    \vspace*{4.2cm}

    % --- 1. Nombre de la actividad ---
    {\sffamily\small\bfseries\color{portadaacento} ACTIVIDAD\par}
    \vspace{0.3cm}
    {\sffamily\fontsize{26}{31}\selectfont\bfseries\color{portadaazul}
      Seminario de Arquitecturas\\[2pt] de Software:\par}
    \vspace{0.4cm}
    {\sffamily\Large\color{portadagris}
      Un Análisis Comparativo de Estilos\\[2pt] Basados en Capas y Dominio\par}
    \vspace{0.55cm}
    {\color{portadaacento}\rule{3.2cm}{2.5pt}\par}

    \vspace{0.9cm}

    % --- 2. Nombre del equipo ---
    {\sffamily\small\color{portadagris} EQUIPO\par}
    \vspace{0.12cm}
    {\sffamily\huge\bfseries\color{portadaazul} Team Bales\par}

    \vspace{0.8cm}

    % --- 3 y 4. Integrantes y contribución individual ---
    {\sffamily\small\color{portadagris} INTEGRANTES Y PREGUNTAS EN LAS QUE TRABAJÓ\par}
    \vspace{0.3cm}
    {\sffamily\normalsize
      \newcommand{\integrante}[2]{\textbf{##1}\quad{\color{portadaacento}\textbullet}\quad{\color{portadagris}##2}\par\vspace{0.22cm}}%
      \integrante{Heily Rodríguez Rodríguez}{1 y 2}
      \integrante{Enrique González Moreira}{4 y 6}
      \integrante{Ernesto Alejandro Soler Choong}{3 y 5}
      \integrante{Mauricio Brindon Carbo}{Sin contribución}
    }

    \vspace{0.8cm}

    % --- 5. Enlace al repositorio público ---
    {\sffamily\small\color{portadagris} REPOSITORIO PÚBLICO\par}
    \vspace{0.12cm}
    {\sffamily\large\href{https://github.com/Team-Bales/seminario-arq-capas}%
      {\color{portadaazul}https://github.com/Team-Bales/seminario-arq-capas}\par}
  \end{titlepage}
}


% \pregunta{número}{ruta sin .tex}
\newcommand{\pregunta}[2]{%
  \IfFileExists{#2.tex}{\setcounter{chapter}{\numexpr#1-1\relax}\input{#2}}{}%
}

% --- Notas de investigación: OCULTAS momentáneamente ---
% Para volver a mostrarlas, cambiar \notasinvestigacionfalse por \notasinvestigaciontrue.
% Mientras están ocultas, se omite el \input de investigacion/... y el título
% «Notas de investigación» que cada solucion.tex pone antes de ese \input.
\newif\ifnotasinvestigacion
\notasinvestigacionfalse

\begin{document}

\ifnotasinvestigacion\else
  \makeatletter
  \let\inputoriginal\input
  \renewcommand{\input}[1]{\in@{investigacion/}{#1}\ifin@\else\inputoriginal{#1}\fi}
  \let\seccionoriginal\section
  \def\notasinvtitulo{Notas de investigación}
  \newcommand{\seccionestrella}[1]{\def\tituloactual{#1}%
    \ifx\tituloactual\notasinvtitulo\else\seccionoriginal*{#1}\fi}
  \renewcommand{\section}{\@ifstar{\seccionestrella}{\seccionoriginal}}
  \makeatother
\fi

% La portada no lleva número; la numeración (arábiga) empieza en 1 con el índice.
\portada
\setcounter{page}{1}
\tableofcontents

% ============================================================
% introduccion.tex — Introducción del documento
% SemArq1_TeamBales.tex, antes del primer capítulo (P1).
% ============================================================

\chapter*{Introducción}
\addcontentsline{toc}{chapter}{Introducción}

Este documento recoge las respuestas del equipo Team Bales a las seis preguntas del seminario. Cada capítulo responde a una pregunta a partir de la bibliografía orientada y termina con las notas de investigación que la sustentan. Las preguntas se abordan de forma independiente, pero comparten un mismo enfoque: analizar cómo se organizan las dependencias entre las partes de un sistema y qué consecuencias tiene esa organización para su mantenimiento y sus pruebas. Al final se incluyen la bibliografía y la declaración de uso de inteligencia artificial.


% Capítulos de contenido (P1..P6). Cada solucion.tex incluye, al final,
% el \input de su notas.tex de investigación.
\pregunta{1}{preguntas/p1-trade-offs/solucion}
\pregunta{2}{preguntas/p2-simetria-hexagono/solucion}
\pregunta{3}{preguntas/p3-ddd-bounded-context/solucion}
\pregunta{4}{preguntas/p4-inversion-dependencias/solucion}
\pregunta{5}{preguntas/p5-kitty-problem/solucion}
\pregunta{6}{preguntas/p6-ncapas-calidad/solucion}

% ============================================================
% conclusiones.tex — Conclusiones del documento
% SemArq1_TeamBales.tex, después del último capítulo (P6) y antes de
% \printbibliography.
% ============================================================

\chapter*{Conclusiones}
\addcontentsline{toc}{chapter}{Conclusiones}

El análisis de las preguntas lleva a una misma idea: lo que distingue a una buena arquitectura no es el número de capas ni la forma del dibujo, sino que la lógica de negocio no dependa de los detalles técnicos. Cuando se consigue, el sistema se puede probar y cambiar por partes sin que cada modificación se propague a todo el código. Esa independencia, sin embargo, tiene un coste de construcción y mantenimiento \fuente{martin2017}{218}, por lo que la arquitectura adecuada depende del tamaño del sistema y de cuánto se espera que cambie.


\printbibliography[heading=bibintoc, title={Referencias bibliográficas}, env=bibliography]

% ============================================================
% declaracion-ia.tex — Uso de IA
% SemArq1_TeamBales.tex, DESPUÉS de \printbibliography, como exige la
% orientación del seminario (sección 3): epígrafe independiente, al
% final del documento, que indique (1) herramienta y versión,
% (2) propósito y (3) forma de integración del resultado.
% Fuente de trabajo: «Uso de IA (Claude).md».
% ============================================================

\chapter*{Uso de IA}
\addcontentsline{toc}{chapter}{Uso de IA}

Para la elaboración de este seminario, el equipo utilizó dos herramientas de inteligencia artificial. Heily Rodríguez Rodríguez y Ernesto Alejandro Soler Choong usaron \textbf{Claude}, de Anthropic, con el modelo \textbf{Claude Opus 5.5}. Enrique González Moreira usó \textbf{OpenCode} con el modelo \textbf{GLM 5.3}.

Las herramientas se emplearon para buscar en la bibliografía recomendada los capítulos, epígrafes y páginas útiles para cada pregunta, y para comprobar que cada cita remite a la página o sección correcta. También se usaron para redactar borradores de las respuestas, mejorar la formalidad y la claridad de la redacción, y traducir al español ideas y fragmentos de las fuentes, que en su mayoría están en inglés. Por último, sirvieron para ayudar a cumplir los requisitos formales de la orientación, como la portada, el estilo de citas y la correspondencia entre las citas del texto y las referencias, y para generar el código \LaTeX{} del documento, incluidos los diagramas.

Los borradores generados se utilizaron como base de las respuestas, que el equipo revisó, corrigió y reescribió hasta su versión final. Las citas se compararon con los textos originales de la bibliografía. Las traducciones sirvieron para comprender y parafrasear las ideas de los autores, y los fragmentos traducidos que aparecen entre comillas se contrastaron con el texto original. El código \LaTeX{} y los diagramas se revisaron y ajustaron hasta obtener el resultado deseado.

Todo el contenido generado fue revisado, verificado y adaptado por el equipo, quien asume la responsabilidad final del documento.


\end{document}
